\documentclass[11pt,a4paper]{article}
\usepackage[T1]{fontenc}
\usepackage[utf8]{inputenc}
\usepackage{lmodern,amsmath,amssymb,booktabs,array,xcolor,microtype}
\usepackage[a4paper,left=2.3cm,right=2.3cm,top=2.5cm,bottom=2.5cm]{geometry}
\usepackage[colorlinks=true]{hyperref}
\definecolor{ojacblue}{rgb}{0,0,0.8}
\hypersetup{linkcolor=ojacblue,citecolor=ojacblue,urlcolor=ojacblue}
\renewcommand{\baselinestretch}{1.08}
\setlength{\parindent}{0pt}
\setlength{\parskip}{4pt}
\setlength{\emergencystretch}{3em}
\newcommand{\R}{https://github.com/ralfanop/Certifications_Eulerian_polar-curve_computations}
\newcommand{\dir}[1]{\href{\R/tree/main/#1}{\texttt{#1}}}
\newcommand{\eq}[1]{(#1)}
\begin{document}

\begin{center}
{\Large\bfseries Manuscript 12: Lean~4 and Arb certification record}\\[6pt]
{\large Bargmann--Fock Representations of the Eulerian Polar Curve: Polyanalytic Extensions and Lanczos Asymptotics}\\[6pt]
Roberto Alfano, 7~November~2026
\end{center}

\section*{1. Scope.}
This record states, for manuscript~12, which results are kernel-checked in Lean~4~[2], which values are certified in Arb ball arithmetic~[1], how the libraries are rebuilt and audited, and which formalisation tasks remain open. It replaces Appendices~B, C and~D of the article (Lean~4 certification) and the description of the supplementary packages formerly given in its Data Availability statement. Equation, section and result numbers are those of the article that cites this record; in it the subsections of Section~5 are numbered 5.1--5.7. All directories named below are in the repository \url{\R}.

\section*{2. The proportional regime: the library \texttt{UniformLanczos}.}
The Lean~4 formalisation of this section comprises 76 theorems and 48 lemmas in 11 modules. The library \texttt{UniformLanczos}, in the directory \dir{lean/UniformLanczos}, is compiled with the Lean~4 toolchain \texttt{v4.35.0-rc3} and Mathlib commit \texttt{0272bbe1c093}; it contains no \texttt{sorry}, and every main declaration depends only on the axioms \texttt{propext}, \texttt{Classical.choice} and \texttt{Quot.sound}. It certifies the following statements.
\begin{itemize}\setlength{\itemsep}{1pt}\raggedright
\item Lemma~5.3 in full, with $\Lambda'$ in the closed form \eq{92}, including the inverse $I'$ and \eq{110} (module \texttt{Lambda}).
\item Lemma~5.4 in full, with $\Phi$ in the $\theta$-form of \eq{98} (\texttt{Endpoint}, \texttt{EndpointConcrete}).
\item Lemma~5.7(b) (\texttt{endpoint\_concrete}).
\item The cancellations \eq{124}--\eq{125} and the identities \eq{96}--\eq{97} (\texttt{Cancellation}).
\item In Lemma~5.8, the formula \eq{137} for $Q'$, the Lipschitz transfer \eq{139}, the passage from Lipschitz to H\"older bounds, $[-r^*,r^*]\subset\mathcal I_{j,M}$, the arcsine substitution in \eq{140} for Lebesgue integrals without any integrability hypothesis, and the equivalence of \eq{101}--\eq{102} with the endpoint equation (\texttt{ClassQ}, \texttt{MRS}, \texttt{Arcsine}).
\item The scaling \eq{142} and the uniformity argument with the filling step (\texttt{Diagonal}).
\item The first two terms of \eq{143}, unconditionally: the bound $|I'(x)-12x-144x^3/5|\leq520x^5$ for $0\leq x\leq\Lambda'(1)$, obtained from $|\Lambda'(s)-s/12+s^3/720|\leq s^5/25600$ for $0<s\leq2$, and the resulting $|F(\alpha)-\alpha/12+\alpha^2/20|\leq K_F\alpha^3$ (\texttt{Corollary}, \texttt{Expansion}).
\item All seven coefficients of \eq{143}, unconditionally, with the remainder \eq{149} (module \texttt{HigherProfileCoefficients}, 30~October~2026, 39 declarations): $|\Lambda'(s)-\sum_{k=1}^{7}B_{2k}s^{2k-1}/(2k)!|\leq s^{15}/(2.5\cdot10^{12})$ for $0\leq s\leq2$, $\Lambda''\geq1/32$ on $(0,2]$, $|I'(x)-I'_{13}(x)|\leq9\cdot10^{5}x^{15}$ and $|\Phi(a)-\Phi_{14}(a)|\leq(9/2)\cdot10^{5}a^{16}$ for $0\leq a,x\leq1/40$, where $I'_{13}$ extends \eq{144} by $179132166144x^{11}/21896875+3789966163968x^{13}/109484375$ and $\Phi_{14}$ extends \eq{145} by $2775854124a^{14}/765625$; the two reversion identities are exact polynomial identities checked by \texttt{ring}, with remainder polynomials of degrees $154$ and $41$ bounded by linear arithmetic.
\end{itemize}
The Fourier--Laplace analysis of Lemma~5.5 is formalised in the library \texttt{Bridge} (below), whose theorem \texttt{lemma\_5\_5\_phi} proves the bound $|\varphi_M'-I'|\leq C/M$. \texttt{UniformLanczos} is also compiled with that toolchain in the Eulerian--Grothendieck project, where the module \texttt{Joined} inserts this bound into \texttt{endpoint\_concrete}, so that Lemma~5.7(b) holds for the actual $\varphi_M$ (theorem \texttt{endpoint\_actual}), and the module \texttt{Profile} proves Lemma~5.7(a), from the conclusions of Lemma~5.6, and Corollary~5.10.

\section*{3. Probabilistic asymptotics and scalar correlation geometry: the Eulerian--Grothendieck project.}
The Lean~4 formalisation of this section comprises 744 theorems, the $418+286+40$ checked declarations of its three supplements. The supplementary Eulerian--Grothendieck project, dated 29~October~2026, records 418 checked declarations in 52 modules: 383 from the previous certificate and 35 additions. It uses Lean~4.24.0 and Mathlib commit \texttt{f897ebcf72cd16f89ab4577d0c826cd14afaafc7}; it is a project of the repository separate from the library \texttt{UniformLanczos}, which uses the newer toolchain stated above. The seven added modules were compiled; unchanged modules were matched by source and compiled-object hashes. A fresh audit of all 418 declarations found only \texttt{propext}, \texttt{Classical.choice} and \texttt{Quot.sound}, with no custom axiom or proof placeholder. The portable script rebuilds all project modules by default.

A second supplement of the same date adds 19 new modules with 286 checked declarations, compiled one module at a time with the same toolchain and the 52 earlier modules untouched (re-verified by hash); its separate audit again found only the three standard axioms, and no \texttt{sorry}, \texttt{admit}, \texttt{axiom} or \texttt{native\_decide}.

For the actual product-law integral in \eq{187} and its reciprocal, Lean proves expansions at every fixed finite order $R$ and fixed $c>1/12$, with remainder $C_{c,R}M^{-R-1}$ for all integers $M\geq1$ (with a constant $C_{c,R}$ whose existence is certified; the printed constant \eq{229} of Theorem~9.10 is certified for $M\geq4096$). All thirteen displayed rational coefficients are identified with the actual integral coefficients, including the $O(M^{-14})$ remainders. The norm remainder is also uniform for all $c\geq c_0>1/12$, with a constant depending only on $c_0$ and $R$. Thus the probabilistic part \eq{246} of the first special-scale statement is formalised, and the second supplement proves, for Mathlib's \texttt{DirichletCharacter.LFunction} of the quadratic character modulo~$5$, the Dirichlet evaluation $L(1,\chi_5)=2\log\varphi/\sqrt5$ used for \eq{247} (summation by parts and dominated convergence at $s=1$). The project also certifies the exact projector \eq{253}, attainment of its scalar maximum, and the global optimisation and critical limit \eq{274} for the Eulerian phase.

The second supplement formalises, in addition: the Gaussian Fock identification of Theorem~9.2, with the vacuum probability as the squared norm of the projection onto the constants in $L^2$ of the Gaussian measure; the branch $0<c\leq1/12$ of Theorem~9.6; the constants of Lemmas~9.3--9.4 and Theorem~9.5; the majorant \eq{224} and the printed all-order constant for $M\geq4096$; parameter-uniform reciprocal remainders; the Bernoulli and zeta identifications of the cumulants at every order; the arbitrary-phase form of Theorem~9.15 with the constant $2000$ for $0<p\leq1/144$; the sharp Lipschitz constant $2/\sqrt5$ on the disk; the bilinear scalar and vector interpretation with $q_{\mathbb C}(P)=4$ and $\mathcal Q\leq K$ for every admissible complex Grothendieck constant $K$; the exposure of the frame Gram matrix on the elliptope, the endpoint decompositions, and Proposition~10.1. Neither convergence of an infinite $1/M$ expansion nor uniformity at the critical threshold is asserted. The supplement contains reproducible programs, enclosures, sources, declaration types, build records and the axiom audit.

A third supplement, dated 2~November~2026, proves Lemma~9.8 at every order (module \texttt{ExactPolynomialAllOrders}, 40 theorems); the proportional-regime coverage of \texttt{UniformLanczos} includes all seven coefficients of \eq{143}.

\section*{4. Sections 3, 4, 6--8 and Lemma 5.5: the library \texttt{Bridge}.}
The proofs of Sections~3, 4, 6, 7 and~8, of Subsections~9.1--9.2, of Lemma~5.5 and of Appendix~A are checked in Lean~4 in the library \texttt{Bridge} of the Eulerian--Grothendieck project, in the directory \dir{lean/EulerianGrothendieck}, with the same toolchain and Mathlib commit. The library, dated 2~November~2026, has 51 modules and 821 theorems; its audit file \texttt{AuditBridge.lean} finds only the axioms \texttt{propext}, \texttt{Classical.choice} and \texttt{Quot.sound}, so no theorem depends on \texttt{sorry} or on a further axiom. The file \href{\R/blob/main/lean/EulerianGrothendieck/README_Bridge_V173.md}{\texttt{README\_Bridge\_V173.md}} of the same directory pairs the numbered statements and equations of these sections with their Lean theorems, in the numbering of the article. After installing the toolchain named in the file \texttt{lean-toolchain}, the commands \texttt{lake exe cache get}, \texttt{lake build Bridge} and \texttt{lake env lean AuditBridge.lean}, run in that directory, fetch the compiled Mathlib, rebuild the library and repeat the audit.

For Lemma~5.5, the twenty-one modules of \texttt{Bridge/Saddle} prove \eq{115} with explicit first and second derivatives (theorem \texttt{lemma\_5\_5}), the continuity of the second derivative on $[-b,b]$ (\texttt{lemma\_5\_5\_C2}) and the two consequences stated after \eq{115} (\texttt{lemma\_5\_5\_phi}). They depart from the printed proof in one point: the tilt $s$ is held fixed while differentiating in $y=Mx$, and $s=s(x)$ is set afterwards, which is legitimate because the Fourier representation of $w_M$ holds for every real tilt. The library also proves the decimal values displayed in \eq{4}--\eq{6}, \eq{12}, \eq{13} and \eq{31}, and $\eta(1/4)$ in \eq{114}.

\section*{5. Fixed-degree Lanczos coefficients: Theorem 5.1 and Corollary 5.2.}
The Lean~4 formalisation of Theorem~5.1 has two principal formal statements, which cover, respectively, the integral Hankel determinants of actual independent uniform sums and the coefficients of a defining monic orthogonal recurrence. For each fixed positive degree they prove eventual nonvanishing of the required denominator and an existential quantitative bound $C_j/M$. The supplementary proof also verifies the exact moment-polynomial bridge for arbitrary independent uniform sums and moment orders, Gaussian Gram positivity, and the first determinant coefficients. Its archived build logs report successful checking with the toolchain of the Eulerian--Grothendieck project; the complete module names, revision identifier, axiom audit and a portable build script are in the directory \dir{lean/FixedDegreeLanczos}, with the exact Lanczos data and the archived sources and logs. On 2~November~2026 the repository's continuous integration recompiled it from source, together with the proofs of Corollary~5.2 and of the exact values of $\beta_1^{[M]}$, $\beta_2^{[M]}$ and $\kappa_4(S_M)$ in Subsection~9.5. The independently checked rational expressions for degrees $5\leq j\leq16$ are included as exact supplementary data.

\section*{6. Numerical certificates.}
Table~1 of the article compares $\beta_j^{[M]}/M^2$ with $F(j/M)$. The original Gauss--Legendre and Stieltjes computations are retained in \dir{numerics/uniform_lanczos}. The supplement of 29~October~2026 independently certifies all 25 displayed values and errors by Arb intervals: exact uniform-convolution moments determine the monic recurrence through parity Gram matrices or a modified Chebyshev recurrence, while validated integration encloses $F$. Every resulting enclosure lies inside its exact rational rounding cell. The approximate constancy of $M\mathcal E_{j,M}$ suggests a $1/M$ error rate, which is not proved by these computations. The displayed saddlepoint value $\eta(1/4)\approx-0.184391$ is independently certified; this checks that coefficient, not the remainder in \eq{115}. The same Arb/FLINT interval computations certify six displayed scalar decimal values. Exact rational rounding or truncation cells are checked; these numerical certificates are not Lean proofs or an asymptotic error-rate theorem.

\section*{7. Supplementary packages.}
The supplementary package of 29~October~2026, the Eulerian--Grothendieck Lean~4 and numerical update, contains its mathematical report, exact rational computations, sources, declaration inventory, build logs, axiom audit and portable verification script. The proportional-regime library and numerical checks are in \dir{lean/UniformLanczos} and \dir{numerics/uniform_lanczos}; the library \texttt{Bridge} is in \dir{lean/EulerianGrothendieck}; the fixed-degree material of Theorem~5.1 is in \dir{lean/FixedDegreeLanczos}.

\section*{References.}
\begin{enumerate}\setlength{\itemsep}{1pt}
\item[{[1]}] F. Johansson, Arb: efficient arbitrary-precision midpoint-radius interval arithmetic, \emph{IEEE Trans. Comput.} \textbf{66} (2017), no.~8, 1281--1292. \url{https://doi.org/10.1109/TC.2017.2690633}.
\item[{[2]}] L. de Moura and S. Ullrich, The Lean~4 theorem prover and programming language, in \emph{Automated Deduction -- CADE 28}, A. Platzer and G. Sutcliffe (eds.), Lecture Notes in Computer Science \textbf{12699}, Springer, Cham, 2021, 625--635. eISBN 9783030798765. \url{https://doi.org/10.1007/978-3-030-79876-5_37}.
\end{enumerate}

\end{document}
